\documentclass[10pt,a4paper]{amsart}
\usepackage{amsmath,amssymb,amsthm,graphicx}
\usepackage{verbatim}
\newenvironment{reference}{\comment}{\endcomment}
%\newenvironment{reference}{}{}
\newenvironment{slogan}{\comment}{\endcomment}
\newenvironment{history}{\comment}{\endcomment}

% For commutative diagrams we use Xy-pic
\usepackage[all]{xy}

% We use 2cell for 2-commutative diagrams.
\xyoption{2cell}
% Theorem environments.
%
\theoremstyle{plain}
\newtheorem{theorem}[subsection]{Theorem}
\newtheorem{proposition}[subsection]{Proposition}
\newtheorem{lemma}[subsection]{Lemma}

\theoremstyle{definition}
\newtheorem{definition}[subsection]{Definition}
\newtheorem{example}[subsection]{Example}
\newtheorem{exercise}[subsection]{Exercise}
\newtheorem{situation}[subsection]{Situation}

\theoremstyle{remark}
\newtheorem{remark}[subsection]{Remark}
\newtheorem{remarks}[subsection]{Remarks}

\numberwithin{equation}{subsection}

% Macros
%
\def\lim{\mathop{\mathrm{lim}}\nolimits}
\def\colim{\mathop{\mathrm{colim}}\nolimits}
\def\Spec{\mathop{\mathrm{Spec}}}
\def\Hom{\mathop{\mathrm{Hom}}\nolimits}
\def\Ext{\mathop{\mathrm{Ext}}\nolimits}
\def\SheafHom{\mathop{\mathcal{H}\!\mathit{om}}\nolimits}
\def\SheafExt{\mathop{\mathcal{E}\!\mathit{xt}}\nolimits}
\def\Sch{\mathit{Sch}}
\def\Mor{\mathop{\mathrm{Mor}}\nolimits}
\def\Ob{\mathop{\mathrm{Ob}}\nolimits}
\def\Sh{\mathop{\mathit{Sh}}\nolimits}
\def\NL{\mathop{N\!L}\nolimits}
\def\CH{\mathop{\mathrm{CH}}\nolimits}
\def\proetale{{pro\text{-}\acute{e}tale}}
\def\etale{{\acute{e}tale}}
\def\QCoh{\mathit{QCoh}}
\def\Ker{\mathop{\mathrm{Ker}}}
\def\Im{\mathop{\mathrm{Im}}}
\def\Coker{\mathop{\mathrm{Coker}}}
\def\Coim{\mathop{\mathrm{Coim}}}

% Boxtimes
%
\DeclareMathSymbol{\boxtimes}{\mathbin}{AMSa}{"02}

%
% Macros for moduli stacks/spaces
%
\def\QCohstack{\mathcal{QC}\!\mathit{oh}}
\def\Cohstack{\mathcal{C}\!\mathit{oh}}
\def\Spacesstack{\mathcal{S}\!\mathit{paces}}
\def\Quotfunctor{\mathrm{Quot}}
\def\Hilbfunctor{\mathrm{Hilb}}
\def\Curvesstack{\mathcal{C}\!\mathit{urves}}
\def\Polarizedstack{\mathcal{P}\!\mathit{olarized}}
\def\Complexesstack{\mathcal{C}\!\mathit{omplexes}}
% \Pic is the operator that assigns to X its picard group, usage \Pic(X)
% \Picardstack_{X/B} denotes the Picard stack of X over B
% \Picardfunctor_{X/B} denotes the Picard functor of X over B
\def\Pic{\mathop{\mathrm{Pic}}\nolimits}
\def\Picardstack{\mathcal{P}\!\mathit{ic}}
\def\Picardfunctor{\mathrm{Pic}}
\def\Deformationcategory{\mathcal{D}\!\mathit{ef}}
\usepackage{cleveref}
\crefname{originallemma}{original lemma}{original lemmas}
\Crefname{originallemma}{Original lemma}{Original lemmas}
\makeatletter
\newlabel{morphisms-lemma-base-change-smooth}{{01VB}{source}{}{sourcepointer-1}{}}
\newlabel{morphisms-lemma-base-change-smooth@cref}{{[originallemma][0][]01VB}{[1][0][]source}{}{}{}}
\makeatother
\crefname{originallemma}{original lemma}{original lemmas}
\Crefname{originallemma}{Original lemma}{Original lemmas}
\makeatletter
\newlabel{morphisms-lemma-composition-proper}{{01W3}{source}{}{sourcepointer-2}{}}
\newlabel{morphisms-lemma-composition-proper@cref}{{[originallemma][0][]01W3}{[1][0][]source}{}{}{}}
\makeatother
\crefname{originallemma}{original lemma}{original lemmas}
\Crefname{originallemma}{Original lemma}{Original lemmas}
\makeatletter
\newlabel{morphisms-lemma-finite-proper}{{01WN}{source}{}{sourcepointer-3}{}}
\newlabel{morphisms-lemma-finite-proper@cref}{{[originallemma][0][]01WN}{[1][0][]source}{}{}{}}
\makeatother
\crefname{originallemma}{original lemma}{original lemmas}
\Crefname{originallemma}{Original lemma}{Original lemmas}
\makeatletter
\newlabel{properties-lemma-normal-dimension-1-regular}{{0BX2}{source}{}{sourcepointer-4}{}}
\newlabel{properties-lemma-normal-dimension-1-regular@cref}{{[originallemma][0][]0BX2}{[1][0][]source}{}{}{}}
\makeatother
\crefname{originallemma}{original lemma}{original lemmas}
\Crefname{originallemma}{Original lemma}{Original lemmas}
\makeatletter
\newlabel{varieties-lemma-dim-1-proper-projective}{{0A26}{source}{}{sourcepointer-5}{}}
\newlabel{varieties-lemma-dim-1-proper-projective@cref}{{[originallemma][0][]0A26}{[1][0][]source}{}{}{}}
\makeatother
\crefname{originallemma}{original lemma}{original lemmas}
\Crefname{originallemma}{Original lemma}{Original lemmas}
\makeatletter
\newlabel{varieties-lemma-finite-extension-geometrically-normal}{{0BXT}{source}{}{sourcepointer-6}{}}
\newlabel{varieties-lemma-finite-extension-geometrically-normal@cref}{{[originallemma][0][]0BXT}{[1][0][]source}{}{}{}}
\makeatother
\crefname{originallemma}{original lemma}{original lemmas}
\Crefname{originallemma}{Original lemma}{Original lemmas}
\makeatletter
\newlabel{varieties-lemma-geometrically-integral}{{038K}{source}{}{sourcepointer-7}{}}
\newlabel{varieties-lemma-geometrically-integral@cref}{{[originallemma][0][]038K}{[1][0][]source}{}{}{}}
\makeatother
\crefname{originallemma}{original lemma}{original lemmas}
\Crefname{originallemma}{Original lemma}{Original lemmas}
\makeatletter
\newlabel{varieties-lemma-geometrically-regular-smooth}{{038X}{source}{}{sourcepointer-8}{}}
\newlabel{varieties-lemma-geometrically-regular-smooth@cref}{{[originallemma][0][]038X}{[1][0][]source}{}{}{}}
\makeatother
\crefname{originallemma}{original lemma}{original lemmas}
\Crefname{originallemma}{Original lemma}{Original lemmas}
\makeatletter
\newlabel{varieties-lemma-normalization-locally-algebraic}{{0BXR}{source}{}{sourcepointer-9}{}}
\newlabel{varieties-lemma-normalization-locally-algebraic@cref}{{[originallemma][0][]0BXR}{[1][0][]source}{}{}{}}
\makeatother
\crefname{originallemma}{original lemma}{original lemmas}
\Crefname{originallemma}{Original lemma}{Original lemmas}
\makeatletter
\newlabel{varieties-lemma-regular-functions-proper-variety}{{04L2}{source}{}{sourcepointer-10}{}}
\newlabel{varieties-lemma-regular-functions-proper-variety@cref}{{[originallemma][0][]04L2}{[1][0][]source}{}{}{}}
\makeatother
\title{Functorial properties of nonsingular projective models under field extension}
\author{The Stacks Project Authors}
\date{stacks-project:curves.tex:local-snapshot-20261012; GNU FDL 1.2 or later}
\begin{document}
\maketitle
\noindent\textit{Editorial note.} All delivered mathematical text is the author's own, copied from the pinned Stacks Project chapter file curves.tex; no statement or proof step was rewritten, reordered or removed.\par\medskip
\noindent\textit{Editorial note.} The counted proof is the author's complete argument: the birational reduction with the author's marked detail, the normalization-properness argument for (1), the purely inseparable normal-regular-smooth chain for (2), and the global-section and base-change arguments for (3) and (4) are all retained verbatim.\par\medskip
\noindent\textit{Editorial note.} The label lemma-smooth-models names the counted statement, and the proof environment immediately following it in the source is delivered verbatim as its proof.\par\medskip
\noindent\textit{Editorial note.} The counted statement is carried by the author's own lemma environment in the chapter source.\par\medskip
\noindent\textit{Allowed standard prerequisites (editorial).}\begin{itemize}
\item Normalizations of proper schemes. Varieties, Lemma varieties-lemma-normalization-locally-algebraic; Morphisms, Lemmas morphisms-lemma-finite-proper and morphisms-lemma-composition-proper; Varieties, Lemma varieties-lemma-dim-1-proper-projective. Makes the normalization of a proper curve finite hence proper, normal, irreducible and birational to the reduced fibre, identifying the smooth model.
\item Improving models by purely inseparable change of fields. Varieties, Lemma varieties-lemma-finite-extension-geometrically-normal; Properties, Lemma properties-lemma-normal-dimension-1-regular; Varieties, Lemma varieties-lemma-geometrically-regular-smooth. Finds a finite purely inseparable K/k making the model geometrically normal, hence geometrically regular for curves, hence smooth.
\item Global functions and base change. Varieties, Lemmas varieties-lemma-geometrically-integral and varieties-lemma-regular-functions-proper-variety; Morphisms, Lemma morphisms-lemma-base-change-smooth. Gives H\textasciicircum{}0 = K for geometrically reduced proper varieties and smoothness of base changes of smooth morphisms.
\end{itemize}\medskip
\ifdefined\hypertarget\hypertarget{sourcepointer-1}{}\fi
\noindent\textit{Source reference.} Original lemma 01VB: sources/stacks-chow/tags.txt, source lines 2440--2440. Prerequisite pointer invoked at proof line(s) 538; it names a standard fact of a sibling chapter that the counted argument uses, and every step of the counted proof built on it is retained verbatim in the delivered body.\par\medskip
\ifdefined\hypertarget\hypertarget{sourcepointer-2}{}\fi
\noindent\textit{Source reference.} Original lemma 01W3: sources/stacks-chow/tags.txt, source lines 2468--2468. Prerequisite pointer invoked at proof line(s) 510; it names a standard fact of a sibling chapter that the counted argument uses, and every step of the counted proof built on it is retained verbatim in the delivered body.\par\medskip
\ifdefined\hypertarget\hypertarget{sourcepointer-3}{}\fi
\noindent\textit{Source reference.} Original lemma 01WN: sources/stacks-chow/tags.txt, source lines 2488--2488. Prerequisite pointer invoked at proof line(s) 509; it names a standard fact of a sibling chapter that the counted argument uses, and every step of the counted proof built on it is retained verbatim in the delivered body.\par\medskip
\ifdefined\hypertarget\hypertarget{sourcepointer-4}{}\fi
\noindent\textit{Source reference.} Original lemma 0BX2: sources/stacks-chow/tags.txt, source lines 14965--14965. Prerequisite pointer invoked at proof line(s) 522; it names a standard fact of a sibling chapter that the counted argument uses, and every step of the counted proof built on it is retained verbatim in the delivered body.\par\medskip
\ifdefined\hypertarget\hypertarget{sourcepointer-5}{}\fi
\noindent\textit{Source reference.} Original lemma 0A26: sources/stacks-chow/tags.txt, source lines 12694--12694. Prerequisite pointer invoked at proof line(s) 514; it names a standard fact of a sibling chapter that the counted argument uses, and every step of the counted proof built on it is retained verbatim in the delivered body.\par\medskip
\ifdefined\hypertarget\hypertarget{sourcepointer-6}{}\fi
\noindent\textit{Source reference.} Original lemma 0BXT: sources/stacks-chow/tags.txt, source lines 14991--14991. Prerequisite pointer invoked at proof line(s) 518; it names a standard fact of a sibling chapter that the counted argument uses, and every step of the counted proof built on it is retained verbatim in the delivered body.\par\medskip
\ifdefined\hypertarget\hypertarget{sourcepointer-7}{}\fi
\noindent\textit{Source reference.} Original lemma 038K: sources/stacks-chow/tags.txt, source lines 4239--4239. Prerequisite pointer invoked at proof line(s) 528; it names a standard fact of a sibling chapter that the counted argument uses, and every step of the counted proof built on it is retained verbatim in the delivered body.\par\medskip
\ifdefined\hypertarget\hypertarget{sourcepointer-8}{}\fi
\noindent\textit{Source reference.} Original lemma 038X: sources/stacks-chow/tags.txt, source lines 4252--4252. Prerequisite pointer invoked at proof line(s) 524; it names a standard fact of a sibling chapter that the counted argument uses, and every step of the counted proof built on it is retained verbatim in the delivered body.\par\medskip
\ifdefined\hypertarget\hypertarget{sourcepointer-9}{}\fi
\noindent\textit{Source reference.} Original lemma 0BXR: sources/stacks-chow/tags.txt, source lines 14989--14989. Prerequisite pointer invoked at proof line(s) 508; it names a standard fact of a sibling chapter that the counted argument uses, and every step of the counted proof built on it is retained verbatim in the delivered body.\par\medskip
\ifdefined\hypertarget\hypertarget{sourcepointer-10}{}\fi
\noindent\textit{Source reference.} Original lemma 04L2: sources/stacks-chow/tags.txt, source lines 5993--5993. Prerequisite pointer invoked at proof line(s) 530; it names a standard fact of a sibling chapter that the counted argument uses, and every step of the counted proof built on it is retained verbatim in the delivered body.\par\medskip
\begin{lemma}
\label{lemma-smooth-models}
Let $k$ be a field. Let $X$ be a geometrically irreducible curve over $k$.
For a field extension $K/k$ denote $Y_K$ a nonsingular projective model
of $(X_K)_{red}$.
\begin{enumerate}
\item If $X$ is proper, then $Y_K$ is the normalization of $X_K$.
\item There exists $K/k$ finite purely inseparable such that $Y_K$ is smooth.
\item Whenever $Y_K$ is smooth\footnote{Or even geometrically reduced.}
we have $H^0(Y_K, \mathcal{O}_{Y_K}) = K$.
\item Given a commutative diagram
$$
\xymatrix{
\Omega & K' \ar[l] \\
K \ar[u] & k \ar[l] \ar[u]
}
$$
of fields such that $Y_K$ and $Y_{K'}$ are smooth, then
$Y_\Omega = (Y_K)_\Omega = (Y_{K'})_\Omega$.
\end{enumerate}
\end{lemma}

\begin{proof}
Let $X'$ be a nonsingular projective model of $X$. Then $X'$ and
$X$ have isomorphic nonempty open subschemes. In particular
$X'$ is geometrically irreducible as $X$ is (some details omitted).
Thus we may assume that $X$ is projective.

\medskip\noindent
Assume $X$ is proper. Then $X_K$ is proper and hence the normalization
$(X_K)^\nu$ is proper as a scheme finite over a proper scheme
(Varieties, Lemma \ref{varieties-lemma-normalization-locally-algebraic}
and Morphisms, Lemmas \ref{morphisms-lemma-finite-proper} and
\ref{morphisms-lemma-composition-proper}).
On the other hand, $X_K$ is irreducible as $X$ is geometrically
irreducible. Hence $X_K^\nu$ is proper, normal, irreducible, and birational
to $(X_K)_{red}$. This proves (1) because a proper curve is projective
(Varieties, Lemma \ref{varieties-lemma-dim-1-proper-projective}).

\medskip\noindent
Proof of (2). As $X$ is proper and we have (1), we can apply
Varieties, Lemma \ref{varieties-lemma-finite-extension-geometrically-normal}
to find $K/k$ finite purely inseparable such that
$Y_K$ is geometrically normal. Then $Y_K$ is geometrically regular
as normal and regular are the same for curves
(Properties, Lemma \ref{properties-lemma-normal-dimension-1-regular}).
Then $Y$ is a smooth variety by
Varieties, Lemma \ref{varieties-lemma-geometrically-regular-smooth}.

\medskip\noindent
If $Y_K$ is geometrically reduced, then $Y_K$ is geometrically
integral (Varieties, Lemma \ref{varieties-lemma-geometrically-integral})
and we see that $H^0(Y_K, \mathcal{O}_{Y_K}) = K$ by
Varieties, Lemma \ref{varieties-lemma-regular-functions-proper-variety}.
This proves (3) because a smooth variety is geometrically reduced
(even geometrically regular, see
Varieties, Lemma \ref{varieties-lemma-geometrically-regular-smooth}).

\medskip\noindent
If $Y_K$ is smooth, then for every extension $\Omega/K$
the base change $(Y_K)_\Omega$ is smooth over $\Omega$
(Morphisms, Lemma \ref{morphisms-lemma-base-change-smooth}).
Hence it is clear that $Y_\Omega = (Y_K)_\Omega$. This proves (4).
\end{proof}

\end{document}
